\begin{table*}[t]
\centering
\small
\caption{Counter-speech generated by Gemma-4-E4B for MultiCONAN messages where evidence cards helped. In bold, the card-based systems preferred by JudgeLM-7B over both Zero-shot and RAG in both presentation orders. EDOS label: predicted by the EDOS classifier and used by the + pred. label systems (MultiCONAN has no EDOS annotation).}
\label{tab:examples-multiconan}
\begin{tabularx}{\textwidth}{@{}p{0.2\textwidth}X@{}}
\toprule
\textbf{Hate speech} & Women are subhuman. \\
\textbf{EDOS label (pred.)} & 2.3 dehumanising attacks \& overt sexual objectification \\
\midrule
\textbf{System} & \textbf{Generated counter-speech} \\
Zero-shot & All human beings, regardless of gender, possess inherent dignity and worth. We must treat everyone with respect and recognize the equal value of every person. \\
RAG & Statements like this reflect sexist assumptions rooted in unequal power relations, which are recognized as a form of discrimination and a violation of human rights. Tackling sexism is a fundamental obligation to guarantee gender equality and prevent violence against women in all spheres of life. \\
RAG + pred. label & Reducing women to subhuman entities is a form of dehumanizing attack that perpetuates gender inequality and is recognized as a form of violence against women and girls. Such sexist hate speech undermines human rights and contributes to a climate of fear and discrimination that limits women's freedom and opportunities. \\
\textbf{Card-RAG} & \textbf{Statements like this are rooted in sexism, which is a form of discrimination that violates fundamental human rights and impedes the full advancement of women in society. We must recognize that violence and hate speech against women are not individual issues, but manifestations of unequal power relations that cause real and severe harm.} \\
\textbf{Card-RAG + pred. label} & \textbf{Reducing women to subhuman entities is a form of dehumanizing attack that violates their dignity and rights, as sexist hate speech aims to objectify and devalue women. We must reject this harmful idea because it perpetuates gender inequality and undermines the fundamental freedom and worth of all women.} \\
Hybrid RAG & Attributing subhuman status to any gender ignores the reality that sexism, rooted in unequal power relations, is a form of discrimination and a violation of human rights. Recognizing this gender inequality is crucial because violence against women is not an individual issue, but a systemic manifestation of sexism. \\
Hybrid + pred. label & Reducing women to subhuman status is a form of dehumanizing attack, which is a recognized category of gender-based violence that aims to humiliate and objectify. Such hateful speech actively undermines women's fundamental human rights, freedom of speech, and dignity. \\
\midrule[\heavyrulewidth]
\textbf{Hate speech} & why do we never talk about female-on-male rape? Isn't this a crime as well? \\
\textbf{EDOS label (pred.)} & 4.2 supporting systemic discrimination against women as a group \\
\midrule
\textbf{System} & \textbf{Generated counter-speech} \\
Zero-shot & All sexual violence is a serious crime, regardless of the gender of the victim or perpetrator. We must advocate for comprehensive safety and justice for everyone affected by sexual assault. \\
RAG & While the focus of some data is on gender-based violence against women, it's important to recognize that legal definitions and societal perceptions of what constitutes a crime vary across countries. These differences in legislation and awareness can significantly impact how violence is reported and understood by the public and authorities. \\
RAG + pred. label & Violence against women, including sexual violence, is a manifestation of unequal power relations and a violation of human rights, not merely an 'individual conflict.' We must recognize that sexism and gender-based violence are systemic issues that require support acknowledging gender inequality. \\
\textbf{Card-RAG} & \textbf{While all sexual violence is a serious crime, the focus on gendered violence is crucial because it highlights how systemic power imbalances between men and women fuel discrimination and human rights violations. Recognizing this gendered context is essential for providing the specific support survivors need.} \\
\textbf{Card-RAG + pred. label} & \textbf{While all sexual violence is a crime, framing violence solely through the lens of gender ignores that violence against women is a manifestation of unequal power relations and systemic discrimination, not just individual conflict. A gendered understanding is crucial because sexism operates at individual, institutional, and structural levels, impacting women's rights and safety across society.} \\
Hybrid RAG & While all sexual violence is wrong, the focus on gender-based violence highlights that sexist hate speech and sexual threats disproportionately target women, and underreporting is a major issue across all forms of violence. We must address the systemic reasons why victims, regardless of gender, fear reporting due to stigma and potential blame. \\
Hybrid + pred. label & Violence against women, including sexual violence, is a manifestation of unequal power relations and a human rights violation, not simply an individual conflict. Recognizing this gendered reality is crucial because sexism and discrimination against women are widespread across all societal levels. \\
\midrule[\heavyrulewidth]
\textbf{Hate speech} & A woman's place is with her husband or his family. This is what they are made for. \\
\textbf{EDOS label (pred.)} & 3.2 immutable gender differences and gender stereotypes \\
\midrule
\textbf{System} & \textbf{Generated counter-speech} \\
Zero-shot & People are capable of so much more than prescribed roles. Every individual deserves the freedom to define their own path and pursue their passions. \\
RAG & Societal expectations about gender roles are shifting, yet sexist behavior persists, as evidenced by the fact that women still perform significantly more unpaid housework than men. True equality requires addressing sexism in the private sphere, aligning with international standards that call for equal rights within the family. \\
RAG + pred. label & Beliefs that dictate a woman's role solely within her husband's family perpetuate harmful gender stereotypes that limit her capacity to make choices and pursue her life plans. True equality recognizes that women have the fundamental right to self-determination across all spheres of life. \\
\textbf{Card-RAG} & \textbf{Defining a woman's worth by passivity or obedience ignores that asserting equal rights is not 'man-hating,' but a fundamental human right. Women deserve the autonomy to make their own decisions about their lives, free from rigid control by family or tradition.} \\
\textbf{Card-RAG + pred. label} & \textbf{Beliefs that dictate a woman's role based on immutable gender differences limit her ability to pursue careers and make choices about her life, as gender stereotypes perpetuate discrimination and restrict potential. Women are capable of diverse roles and contributions, and their value is not defined by their relationship to a husband or family.} \\
Hybrid RAG & Societal expectations often confine women to traditional roles, yet evidence shows that women continue to perform significantly more unpaid domestic labor than men, while asserting equal rights is frequently met with labels like 'unwomanly.' True partnership values individual autonomy and shared responsibility, not rigid control over a woman's life or choices. \\
Hybrid + pred. label & The idea that a woman's role is strictly defined by her husband or family ignores that gender stereotypes limit everyone's ability to pursue their full potential and make their own life choices. Women, like all individuals, deserve the freedom to define their own roles and live without being constrained by rigid, sexist expectations. \\
\bottomrule
\end{tabularx}
\end{table*}
